\section{Introduction}
\label{sec:introduction}

Rotating machinery -- including bearings, gearboxes, rotors, and shafts -- forms the backbone of modern industrial systems ranging from wind turbines and aerospace engines to manufacturing equipment and railway traction systems. Unexpected failure of these critical components can cause catastrophic downtime, with estimated costs reaching millions of dollars per hour in continuous production environments~\cite{hengel2014survey,raihani2017condition}. Accurate and timely fault diagnosis is therefore essential for ensuring operational safety, reducing maintenance costs, and enabling predictive maintenance strategies in Industry 4.0~\cite{lei2020machinery}.

Deep learning has emerged as a dominant paradigm for vibration-based fault diagnosis, demonstrating remarkable success in learning hierarchical features from raw sensor data~\cite{wang2019deep,zhang2020convolutional}. Convolutional Neural Networks (CNNs) and their variants have been widely adopted for processing both temporal waveforms and time-frequency representations~\cite{inoue2018cnn,wen2018new}. More recently, Vision Transformers (ViTs) have been applied to fault diagnosis tasks, leveraging self-attention mechanisms to capture long-range dependencies in vibration signals~\cite{ding2022transformer}. Despite these advances, three fundamental challenges persist under real-world industrial conditions:

First, vibration signals exhibit \textit{time-frequency heterogeneity} across varying operating conditions (load, speed, and environmental factors). Transient impacts from fault-induced impulses coexist with steady-state harmonic components, and their characteristics shift substantially as operating conditions change. Fixed-window spectral analysis or single-domain feature extraction cannot simultaneously capture both transient dynamics and frequency patterns, leading to incomplete representations that degrade under domain shifts. Second, \textit{severe data scarcity} is inherent in industrial settings. Labeled fault samples are expensive and hazardous to collect, as they require running machinery to failure under controlled conditions. The long-tail distribution of natural fault occurrences further compounds this imbalance, with normal operation samples vastly outnumbering fault instances. Third, \textit{heavy background noise} from other machinery, electromagnetic interference, and structural vibrations masks incipient fault signatures in real-world environments, causing conventional methods to degrade sharply when test conditions deviate from clean training data.

To address these challenges, three lines of research have been pursued. Physics-informed methods~\cite{wang2020physics,shen2023physics} incorporate domain knowledge through physics-guided loss functions or wavelet-based convolutional layers, improving generalization but often at the cost of increased architectural complexity. Few-shot meta-learning approaches~\cite{li2021meta,sun2022multi} leverage episodic training paradigms to enable rapid adaptation from limited labeled data, yet existing meta-learning frameworks lack physical consistency guarantees and are vulnerable to overfitting when domain shifts are severe. Dual-stream and multi-modal architectures~\cite{zhao2021fdfnet,han2022mcformer} have been explored for time-frequency fusion, but their fusion mechanisms typically rely on simple concatenation or weighted averaging, lacking physically meaningful alignment between temporal and spectral representations.

In this paper, we propose TAFN (Time-Frequency Aware Fusion Network), which jointly addresses all three challenges through a principled integration of dual-stream representation learning, physics-guided feature fusion, and entropy-regularized meta-learning. Our contributions are as follows:

\begin{enumerate}
    \item \textbf{Dual-stream time-frequency architecture.} We propose a novel dual-stream network that processes raw temporal waveforms through dilated 1D convolutions and Mel-scaled spectrograms through lightweight 2D CNNs with channel attention in parallel. A cross-modal attention mechanism then aligns these heterogeneous representations, capturing both transient impact dynamics and steady-state spectral patterns.
    \item \textbf{Physics-guided fusion with statistical priors.} We embed dimensionless statistical metrics (variance, skewness, kurtosis, and clearance factor) as physically meaningful inductive biases into the fusion process. These metrics, derived from classical vibration analysis, serve as noise-robust anchors that remain stable under varying operating conditions and guide the network toward physically consistent representations.
    \item \textbf{Entropy-regularized meta-learning.} We introduce predictive entropy minimization into the prototypical network training objective. This regularization term encourages confident and well-separated feature embeddings during few-shot adaptation, preventing overfitting to spurious correlations in limited data and improving generalization to unseen operating conditions.
    \item \textbf{State-of-the-art experimental results.} Comprehensive experiments demonstrate that TAFN achieves 92.1\% accuracy with only 1.5M parameters, outperforming ViT-B/16 (88.4\%) while using 57$\times$ fewer parameters. Under extreme noise (-8dB SNR), TAFN maintains 61.8\% accuracy -- a 16.6\% absolute improvement over CNN baselines. In few-shot settings with only 5 samples per class, TAFN reaches 80.5\%, surpassing MAML (68.5\%) and Prototypical Networks (71.3\%).
\end{enumerate}

The remainder of this paper is organized as follows. Section~\ref{sec:related} reviews related work. Section~\ref{sec:problem} formulates the problem. Section~\ref{sec:method} details the TAFN architecture. Section~\ref{sec:experiments} presents experimental results and analysis. Section~\ref{sec:ablation} provides ablation studies. Section~\ref{sec:discussion} discusses implications and limitations. Section~\ref{sec:conclusion} concludes the paper.
